\section{Limitations and threats to validity}
\label{sec:limits}

\begin{itemize}
\item \textbf{Interim safety run.} D1 has one seed and four of five configurations; the Q4\_K\_M Qwen
  reviewer could not start within the VRAM budget beside the owner's server. Intervals are reported and
  hypothesis tests are deferred to the complete multi-seed run.
\item \textbf{Subsample and floor.} D1 uses the first five user tasks and first three injection tasks of
  each AgentDojo suite in registry order, not a random sample, and attack success without any defence was
  near zero, so the run cannot discriminate between reviewers on security.
\item \textbf{Decoding in the earlier run.} The earlier D1 run was configured for temperature 0, but
  AgentDojo's client omits a zero temperature, so its agent sampled at the server's defaults (equal to
  Jig's) without a seed, while its reviewers ran at temperature 0. The protocol v0.4 run sends Jig's
  defaults and the seed explicitly on every request and records them per trial.
\item \textbf{Real DNS in a simulated benchmark.} Jig's \texttt{no-local-network} rule resolves every URL's
  host in real DNS. A D1 run therefore makes DNS lookups for AgentDojo's host names (and no other network
  traffic, because AgentDojo's tools are simulated). Fictional hosts that do not resolve are refused in
  every condition, which lowers Slack utility for reasons unrelated to the reviewers.
\item \textbf{Per-layer attribution is partial.} Which layer refused each call is recorded only for trials
  after a harness change during the run: none of the Granite Guardian trials and 87 of 140 no-reviewer
  trials (in which every refusal is necessarily a core rule).
\item \textbf{Guardian criterion.} Granite Guardian is driven with a defensive criterion written for this
  programme; a different criterion could trade utility for coverage differently.
\item \textbf{Concurrent development.} The runs used Jig's working tree while other work on Jig was in
  progress; each run session records the commit and a hash of the uncommitted core changes.
\item \textbf{Memory pilot scale.} The C1 pilot has one question per type and four sessions per question,
  so intervals are very wide and no difference reported here is statistically established.
\item \textbf{Output limit and hidden reasoning.} The agent's hidden reasoning counts towards its
  8192-token output limit. On every rolling-summary attempt in the full run so far (5 attempts on 3
  questions), the summary rewrite reached the limit with no visible output. Diagnostic replays show the
  model spending the whole budget counting the summary's words one by one to honour the 400-word bound
  (27{,}000--29{,}000 characters of reasoning, not a repetition loop). Those trials are harness errors and are not scored. From protocol
  v0.4 the summary prompt sets no word limit and these trials are re-run as a separate run. The cut-off
  trials are kept as a finding: with a reasoning model, an explicit length limit in the prompt can use up
  the whole output budget on self-checking.
\item \textbf{One agent.} The agent is one ternary 27B model; its Q4 counterpart is downloaded and queued.
\item \textbf{Judge.} The judge is a local model, not the GPT-4o judge LongMemEval's prompts were
  validated with; a hand audit of a sample of verdicts is planned for the full run.
\item \textbf{Simulated time.} Jig's system prompt states the real date, which can conflict with the
  date of a question; memory timestamps are real, so dates are carried in the memory text.
\item \textbf{Shortened histories.} To fit the compute budget we replay a subsample of each question's
  sessions (the answer-bearing ones plus distractors), which makes retrieval easier than in the full
  LongMemEval-S haystack.
\item \textbf{Coverage of the literature.} The review used arXiv only.
\end{itemize}
